\subsection{Large-minimum even terminal packets at a root (1,2,v)}\label{proof:Q5-C180}
\paragraph{C180.0 Fixed directions and statement}



The internal root residual is \((1,\allowbreak \allowbreak  2,\allowbreak \allowbreak  v)\), \(v\ge3\), and the terminal residual satisfies
\[
 4\le a_1<b_1<\cdots<a_j<b_j,\qquad j\ge1,              \tag{C180.1}
\]
The finite positive pair count \(j\), lengths and length gaps are unbounded. The original \(L,\allowbreak \allowbreak  M\ge1\) are independent. Genuine equal pairs may be attached at both roots.

% Replace a pure model-definition block only; keep local root/terminal data.
Use the original finite-space model and arm messages of
Section~\ref{sec:odd-common-calculus}.
When the original \(K\)-notation is used, \(K=-B\) and
\(g=K\one=-u\); in either notation
\(\delta_j=\|B^j\one\|_2^2=\|K^j\one\|_2^2\).
All means, norms, adjoints and spectral integrals use the same
original probability measure.
Fix the actual terminal and root pair as
\[
\begin{gathered}
F=\prod_{i=1}^j h_{a_i}f_{b_i},\qquad P=\bar f_2W_v,\\
\qquad H=pP T^LF.
\end{gathered}\tag{C180.2}
\]
\phantomsection\label{part:C180:dirs}The four arm types have their specified original directions. The unit root arm points outward and gives \(p\); inner edges point from terminal to root and outer edges from root to receiver. The direction is the same for every \(L\) and host.

Set
\[
 \lambda_j=\frac{29}{5}+\frac{64}{5}(j-1).
\]
For every listed integer, legal original kernel and genuine equal-pair decoration,
\[
 \boxed{EH^4\le e^{-\lambda_jd}\le e^{-29d/5}.}          \tag{C180.3}
\]
This fourth moment includes the genuine original unit root arm and complete root pair.

The outer output satisfies
\[
\begin{gathered}
I_q(T^MH)^3\le e^{-[2\lambda_j-c(M)]d}\le e^{-82d/95},\\
\qquad I_q(X)=E[q^{-1/3}|X|^{8/3}],
\end{gathered}\tag{C180.4}
\]
with the all-length Q5-C164 costs
\[
\begin{gathered}
c(1)=204/19,\quad c(2)=10,\quad c(3)=0,\quad\\
c(4)=7,\quad c(M)=6\ (M\ge5).
\end{gathered}\tag{C180.5}
\]

For two such sides on an original five-branch path and the compatible negative zero-step central table, the four positive core lengths remain independent and
\[
 \boxed{Z_D\le e^{-41d/190}\le1.}                     \tag{C180.6}
\]
The stronger exponent is \([2\lambda_-+2\lambda_+-c(M_-)-c(M_+)]/8\).

\paragraph{C180.1 Finite calculus and the actual unequal-pair losses}

Delete only zero-mass atoms. Antisymmetry gives \(K(x,\allowbreak \allowbreak  x)=0\), so \(p,\allowbreak \allowbreak  q\ge\mu(x)>0\). The kernels of \(A,\allowbreak \allowbreak  T\) lie in \([0,\allowbreak \allowbreak  2]\). Constant labels prove strict positivity of every actual tree message. Every norm, mean, spectral measure and Hölder application uses the original \(\mu\).

The finite calculus gives
\[
\begin{gathered}
0\le d\le1/3,\quad\theta=13/32,\quad\|B\|^2\le\theta,\\
\quad |BV(x)|\le\|V\|_2,\quad \delta_{r+1}\le\theta \delta_r.
\end{gathered}
\]
Put \(G_n(z)=\sum_{k=0}^{n-1}z^k\), \(G_0=0\), \(\beta_1=\beta_2=0\), and
\[
 \beta_n=\sum_{r=1}^{\lfloor(n-1)/2\rfloor}(-1)^{r-1}\delta_r.
\]
The exact finite positive convolution is
\[
\begin{gathered}
S_n=1+G_n(B)u=\bar f_n+\sum_{r=3}^n\beta_r \bar f_{n-r},\\
\qquad d-\delta_2\le\beta_r\le d\quad(r\ge3).
\end{gathered}\tag{C180.7}
\]
Full reversal gives the \(f_n\) identity. The four path types lie in \([0,\allowbreak \allowbreak  2]\); \(f_n,\allowbreak \allowbreak  \bar f_n\) have means at most one. The genuine paired products \(f_lh_l,\allowbreak \allowbreak  \bar f_lW_l\) lie in \([0,\allowbreak \allowbreak  1]\).

The spectral measure of \(u\) is supported on \([0,\allowbreak \allowbreak  \theta]\), has mass \(d\), and moments \(\int x^r\,\allowbreak \allowbreak  d\nu=\delta_{r+1}\). At \(z=i\sqrt x\), write \(G_n=E_n+zD_n\). For \(n\ge2\) both coordinates lie in \([1-x,\allowbreak \allowbreak  1]\), while \(n=1\) gives \((1,\allowbreak \allowbreak  0)\). For even \(n\), \(E_n=D_n\); for odd \(n\), \(E_n=1-xD_n\). For \(n\ge3\), \(E_n\le1-x+x^2\), \(|G_n|\le1\), while \(|G_2|^2=1+x\).

For a real polynomial \(Q=E+zD,\allowbreak \) the actual mean and norm are
\[
\begin{gathered}
E[Q(B)u]=-\int D\,d\nu,\qquad\\
\|Q(B)u\|_2^2=\int(E^2+xD^2)d\nu.
\end{gathered}\tag{C180.8}
\]
The global moments are integrated restrictions on the spectral measure.

The actual unequal-pair bound \eqref{eq:L96:long} is
\[
\begin{gathered}
h_a f_b\le\exp\left(\sum_{r=a+1}^bK^r1-16d/5\right)\\
\quad(4\le a<b).
\end{gathered}\tag{C180.9}
\]
Consequently one common finite vector gives
\[
\begin{gathered}
R=\sum_{i=1}^j\sum_{r=a_i+1}^{b_i}K^r1,\qquad\\
\Gamma=(16j/5)d,\qquad F\le e^{R-\Gamma}.
\end{gathered}\tag{C180.10}
\]
Every unequal pair retains its actual scalar loss and original arms.

\paragraph{C180.2 One mask for every finite pair count}

The original exponent intervals are disjoint. Their finite zero-one
mask, with either kernel sign, satisfies
\eqref{eq:odd-common:mask} by Lemma~\ref{lem:odd-common:mask}.

Write \(R=BV_0\). Its lowest original power is five, so \(V_0\) acting on \(u\) starts at \(B^3\), and \(R\) at \(B^4\). Thus
\[
\begin{gathered}
\|V_0\|^2\le(101/50)\theta^3d<d/7,\\
\qquad \sigma:=ER^2\le(101/50)\theta^4d<d/18.
\end{gathered}\tag{C180.12}
\]
The first comparison has numerator \(221897\), denominator \(1638400\), and \(1553279<1638400\) after multiplication by seven. The second has numerator \(2884661\), denominator \(52428800\), and \(51923898<52428800\) after multiplication by eighteen. The row bound gives \(|R|^2<d/7\le1/21<1/16\), hence \(|R|<1/4\). For positive \(R\), consecutive Taylor-tail ratios after the quadratic term are at most \(1/12\), so their coefficient is at most \(6/11<5/8\). For negative \(R\), the integral coefficient is at most \(1/2\). Therefore
\[
 e^R\le1+R+(5/8)R^2.                                \tag{C180.13}
\]

Put \(D_k=EB^kR\). The exact skew moments are \(EK^{2r}\one=(-1)^r\delta_r\), with odd moments zero. Since original powers start at five, the negative \(D_0\) subseries starts at \(\delta_3\), its positive subseries at \(\delta_4\). The positive \(D_1=-EKR\) subseries starts at \(\delta_3\), its negative one at \(\delta_4\). For every fixed \(k\ge2\), the positive subseries starts no earlier than \(\delta_4\); same-sign moment indices differ by two. With \(C_\beta=(1-\theta^2)^{-1}=1024/855\),
\[
\begin{gathered}
-Cbeta\,\delta_3\le D_0\le Cbeta\,\delta_4,\\
\quad -Cbeta\,\delta_4\le D_1\le Cbeta\,\delta_3,\\
\quad D_k\le Cbeta\,\delta_4\quad(k\ge2).
\end{gathered}\tag{C180.14}
\]
This sign classification covers every finite \(k\). Even when the positive-start parity changes, it starts no earlier than \(\delta_4\), with the same two-index spacing. Since \(C_\beta\theta^2=169/855<1/5\), \(C_\beta\theta^3=2197/27360<1/12\),
\[
\begin{gathered}
D_0\in[-d/5,d/12],\quad D_1\in[-d/12,d/5],\quad\\
D_k\le d/12\ (k\ge2).
\end{gathered}\tag{C180.15}
\]
The positive bound for D1 is d/5.

Define
\[
 \delta=D_0+(5/4)\sigma,\qquad\epsilon=\delta+D_1.
\]
The same quantities satisfy
\[
\begin{gathered}
\delta\le11d/72,\quad |\delta|\le d/5,\quad\\
\epsilon\le127d/360,\quad|\epsilon|\le3d/8.
\end{gathered}\tag{C180.16}
\]
For example \(-17d/60\le\epsilon\le127d/360\), both within \([-3d/8,\allowbreak \allowbreak  3d/8]\). Put \(\gamma_k=\beta_k-D_{k-1}\). For every finite \(k\ge3\),
\[
 \gamma_k\ge(19/32-1/12)d=49d/96\ge d/2.              \tag{C180.17}
\]
Only the finite alternating-moment bounds have been used.

\paragraph{C180.3 A root envelope and three complete finite bases}

Let \(n=v-2\ge1\), \(G_n=E+zD\), \(\beta_v=\int E\), \(\beta_{v-1}=\int D\). Add the actual \(\bar f_2\) and \(W_v\), retaining their original negative terms:
\[
 \bar f_2+W_v\le2-\beta_v-\beta_{v-1}-B^2G_n(B)u+\beta_{v-1}u.
\]
For nonnegative \(x,\allowbreak \allowbreak  y\), \(xy\le e^{x+y-2}\). Since \(u\le1\), dropping \(\beta_{v-1}(u-1)\le0\) gives the valid weaker envelope
\[
 P\le\exp\{-\beta_v+B[-BG_n(B)u]\}.                  \tag{C180.18}
\]
Thus root scalar plus mean is
\[
 -\beta_v+E[-B^2G_n(B)u]=-\int(E+xD)d\nu\le-5d/6.   \tag{C180.19}
\]
For odd \(n\), \(E+xD=1\); for even \(n\), \(E=D\ge1-x\), so \(E+xD\ge1-x^2>5/6\). At \(n=1\), the root vector is exactly \(-Bu\), with loss \(d\). Retaining the \(G_2\) exception gives
\[
\|BG_n(B)u\|^2\le\theta(1+\theta)d=(585/1024)d<(16/25)d.\tag{C180.20}
\]

The two short bases below pay their same-frequency root loss and norm jointly.

\paragraph{C180.3.1 The two-step base}

Let \(A_n(z)=1-zG_n(z).\) Then
\[
 |A_n|^2=(1+xD)^2+xE^2\le(1+x)^2+x<(8/5)^2.           \tag{C180.21}
\]
The joint inequality is
\[
 \int\{-4-4E-4xD+(17/4)|A_n|^2\}d\nu\le(17/16)d.  \tag{C180.22}
\]
At \(n=1\), the bracket is \(-15/4+17x/4\). For \(n\ge2\), its derivative in \(E\) is \(-4+17xE/2<0\), and in \(D\) is \(x(9/2+17xD/2)\ge0\). The coordinate rectangle therefore reduces to \(E=1-x,\allowbreak \allowbreak  D=1\), giving
\[
 -15/4+51x/4-17x^2/4+17x^3/4.
\]
This polynomial increases on \([0,\allowbreak \allowbreak  \theta]\), with right endpoint \(132805/131072<17/16\).

\paragraph{C180.3.2 The three-step base}

Let \(B_n(z)=G_2(z)-zG_n(z).\) Its squared modulus is
\[
\begin{gathered}
(1+xD)^2+x(1-E)^2\le(1+x)^2+x^3\\
\le66997/32768<(23/16)^2.
\end{gathered}\tag{C180.23}
\]
The full joint estimate is
\[
 \int\{-8-4E-4xD+(17/4)|B_n|^2\}d\nu\le-3d.        \tag{C180.24}
\]
At \(n=1\), the bracket is \(-31/4\). For \(n\ge2\), it decreases in \(E\) and increases in \(D\), giving the endpoint polynomial
\[
 -31/4+17x/2+17x^2/4+17x^3/4
\]
This increases, with value \(-433915/131072<-3\) at \(\theta\).

\paragraph{C180.3.3 Every longer finite base}

For \(m\ge3,\allowbreak \allowbreak  n\ge1\), write \(G_m=U+zV\). The finite coordinates give \(U+xD\le1+x^2\), \(|V-E|\le x\), including \(n=1\). Hence
\[
\begin{gathered}
|G_m-zG_n|^2\le(1+x^2)^2+x^3\\
\le1493553/1048576<36/25.
\end{gathered}\tag{C180.25}
\]
The final comparison is \(37338825<37748736\), valid for every finite \(m,\allowbreak \allowbreak  n\).

\paragraph{C180.4 Actual positive propagation and the short connectors}

Use the signed full-real-vector estimate
\[
 Ee^{4BV}\le\exp\{4EBV+(17/4)\|V\|^2\}.             \tag{C180.26}
\]
Its mean may have either sign; each application below retains the complete vector.

The terminal envelope, Taylor bound and \(T(R^2)\le2\sigma\) give
\[
 T F\le e^{-\Gamma}(1+u+BR+\delta).
\]
The bracket dominates the actual positive message. Repeated application of the same original positive operator gives
\[
\begin{gathered}
T^LF\le e^{-\Gamma}Y_L,\\
\quad Y_L=\bar f_L+\delta \bar f_{L-1}+B^LR\\
+\sum_{i=1}^{L-1}D_i \bar f_{L-1-i}>0.
\end{gathered}\tag{C180.27}
\]
This is finite Duhamel induction. Combining \(\log p\le-u\), \(\log Y_L\le Y_L-1\) and the root envelope gives the following actual exponential bounds.

\paragraph{C180.4.1 The unit connector}

The full vector is \(V=R-BG_n(B)u,\allowbreak \) and root scalar plus mean gives
\[
 S\le-\Gamma+\epsilon-5d/6\le-\Gamma-173d/360.
\]
The packet and root bounds give \(\|V\|\le21\sqrt d/20\), so
\[
\begin{gathered}
\log EH^4\le-4\Gamma-173d/90+7497d/1600\\
\le-[10+64(j-1)/5]d.
\end{gathered}\tag{C180.28}
\]
At \(j=1\), the exact margin over loss ten is \(527/14400\).

\paragraph{C180.4.2 The two-step connector}

Retaining the actual second, first and zeroth paths gives the full vector
\[
 V=A_n(B)u+BR+\delta\one.
\]
After the jointly paid base is removed, additional scalar plus mean is \(\epsilon+D_2\). Since \(\sqrt d\le3/5\), \(d\le1/3\),
\[
 \|BR\|^2\le(13/576)d<(16/625)d,
 \quad |\delta|\le(3/25)\sqrt d,
\]
Thus \(\|BR+\delta\one\|\le7\sqrt d/25\). This retains the constant coordinate and all cross terms by the triangle inequality. Its norm increment costs
\[
 (17/4)\{2(8/5)(7/25)+(7/25)^2\}d=10353d/2500.
\]
The actual additional scalar cost is \(4\epsilon+4D_2\le157d/90\). Therefore
\[
\begin{gathered}
\log EH^4\le-4\Gamma+\\
\relax [17/16+157/90+10353/2500]d\\
<-4\Gamma+7d=-\lambda_jd.
\end{gathered}\tag{C180.29}
\]
The bracket is less than seven because \(10353\cdot720=7454160<7547500=3019\cdot2500\), and \(7-17/16-157/90=3019/720\). This includes the actual positive \(D_1\) cost \(d/5\).

\paragraph{C180.4.3 The three-step connector}

The exact \(\bar f_3=1-d+u+Bu+B^2u\) gives
\[
 V=B_n(B)u+\delta u+B^2R+\epsilon\one.
\]
Outside the jointly paid base, scalar plus mean is \(\epsilon+D_2+D_3-\delta d\). The packet and \(d\le1/3,\allowbreak \allowbreak  \sqrt d\le3/5\) bound the perturbation norm by
\[
 [1/15+1/10+9/40]\sqrt d<(5/12)\sqrt d.
\]
Here \(\|B^2R\|^2\le\theta^2d/18<d/100\). The base gives norm increment at most \(6715d/1152<6d\), while additional scalar cost is at most
\[
 4\epsilon+4D_2+4D_3-4\delta d\le(211/90)d<(12/5)d.
\]
Thus
\[
\begin{gathered}
\log EH^4\le-4\Gamma-3d+6d+(12/5)d\\
\le-[7+64(j-1)/5]d.
\end{gathered}\tag{C180.30}
\]

\paragraph{C180.5 Every finite length at least four}

Substitution of the exact positive convolution gives
\[
\begin{gathered}
Y_L=S_L+\delta \bar f_{L-1}+D_1\bar f_{L-2}+B^LR\\
-\sum_{k=3}^L\gamma_k \bar f_{L-k}.
\end{gathered}\tag{C180.31}
\]
Both \(k=L,\allowbreak \allowbreak  L-1\) occur in the actual finite sum when \(L\ge4\), giving the original \(\bar f_0,\allowbreak \allowbreak  \bar f_1\). Retain them using the proved gamma signs, drop only other negative terms, and use \(\bar f_n\le2\) on positive-coefficient terms:
\[
 Y_L\le S_L+B^LR+2\delta_++2(D_1)_+-(d/2)(2+u).
\]
The right side dominates positive \(Y_L.\) After multiplying by the actual p and root pair, the full exponent vector is
\[
 V=[G_{L-1}(B)-BG_n(B)]u+B^{L-1}R-(d/2)1.             \tag{C180.32}
\]
The exact \(B\)-mean is \(-\beta_{L+1}+D_L=-\gamma_{L+1}\le-d/2\); the extra endpoint index is covered. Combining with the root mean gives
\[
\begin{gathered}
S\le-\Gamma+2\delta_++2(D_1)_+-3d/2-5d/6\\
\le-\Gamma-293d/180.
\end{gathered}\tag{C180.33}
\]
Indeed \(2\delta_++2(D_1)_+\le11d/36+2d/5=127d/180\).

The packet gives \(\|B^{L-1}R\|^2\le\theta^3d/18<d/256\); the finite base has norm coefficient \(6/5\). The genuine constant \((d/2)\one\) has norm at most \(3\sqrt d/10\). The triangle bound retains all cross terms:
\[
 \|V\|\le[6/5+1/16+3/10]\sqrt d=(25/16)\sqrt d.
\]
The signed full-vector estimate now gives
\[
\begin{gathered}
\log EH^4\le-4\Gamma-293d/45+10625d/1024\\
\le-[8+64(j-1)/5]d.
\end{gathered}\tag{C180.34}
\]
For example \(10625/1024<21/2\), while \(64/5+293/45-21/2>8\). Every arbitrarily large finite connector is covered.

Comparing the four length ranges, the smallest uniform loss is the two-step value \(\lambda_j\). The same direction works on the entire domain and every adjacent or unit-length boundary.
